\documentclass[11pt]{article}
\usepackage[utf8]{inputenc}
\usepackage[T1]{fontenc}
\usepackage[italian]{babel}
\usepackage{geometry}
\geometry{a4paper, margin=1in}
\usepackage{xcolor}
\usepackage{listings}
\usepackage{hyperref}

\usepackage{tcolorbox}
\tcbuselibrary{breakable}
\begin{document}

\begin{tcolorbox}[colback=blue!5!white,colframe=blue!75!black,title=\textbf{DOMANDA}, breakable]
Esempio di web server esteso con gestione CGI
\begin{itemize}
    \item Aprire il file \texttt{serverHTTP-CGI.c} in Esempi-web/ e analizzarne il contenuto
    \item Compilarlo come al solito ed eseguirlo
    \item Aprire il browser preferito e il file \texttt{sommatrice-web.html}
    \begin{itemize}
        \item Cosa si vede sul browser?
        \item NOTA: Provare con numeri positivi, negativi, con parte decimale...
    \end{itemize}
    \item A cosa corrisponde il secondo parametro della funzione \texttt{sommatrice()} ?
\end{itemize}
\end{tcolorbox}

\begin{tcolorbox}[colback=green!5!white,colframe=green!75!black,title=\textbf{RISPOSTA}, breakable]
\textbf{Soluzione Esercizio 27:} Web server con gestione CGI

\textbf{Operazioni Preliminari}\\
Ho copiato i file \texttt{serverHTTP-CGI.c} e \texttt{sommatrice-web.html} nella nuova cartella dell'esercizio e ho compilato l'eseguibile del server C, per poi avviarlo.

\textbf{Risposte alle domande:}

\begin{enumerate}
    \item \textbf{Cosa si vede sul browser aprendo sommatrice-web.html?}\\
    Si presenta una semplice interfaccia (creata tramite Web Form) che richiede l'inserimento di "Operando 1" e "Operando 2", accompagnata dal pulsante "Somma". Inviando il modulo, il browser invia i dati inseriti appendendoli all'URL (es: \texttt{/sommatrice?op1=5.5\&op2=-2.2}).\\
    Provando a inviare numeri positivi, negativi o con la virgola (utilizzando il punto come separatore decimale), il server restituisce correttamente il risultato matematico elaborato (es. Risultato=3.300000).\\
    Questo avviene perché il backend scritto in C utilizza la funzione \texttt{atof()}, la quale supporta nativamente il parsing delle stringhe contenenti segni negativi e punti decimali.

    \item \textbf{A cosa corrisponde il secondo parametro della funzione \texttt{sommatrice()}?}\\
    \textbf{Analizzando la firma della funzione:} \texttt{void sommatrice(char *url, FILE *out)}\\
    Il parametro \texttt{FILE *out} corrisponde allo stream di output del socket \texttt{TCP}, ovvero la connessione aperta con il client (il browser). \\
    Se guardiamo il \texttt{main()}, notiamo che la funzione viene richiamata passando \texttt{connfd} (il file descriptor della connessione): \texttt{sommatrice(url, connfd);}.\\
    \textbf{Questo meccanismo è l'essenza dell'interfaccia CGI (Common Gateway Interface):} il server "delega" l'elaborazione a un programma o funzione esterna, passandole i parametri e fornendole l'accesso diretto allo standard output della rete. Qualsiasi cosa la funzione \texttt{sommatrice} stampi tramite \texttt{fprintf(out, ...)}, andrà a finire dritta al browser (ecco perché la funzione si preoccupa di stampare lei stessa l'intestazione \texttt{HTTP 200 OK} e i tag HTML).
\end{enumerate}

\textbf{Comandi utili:}
\begin{lstlisting}[language=bash, basicstyle=\ttfamily\small]
# Compila ed esegui
gcc serverHTTP-CGI.c network.c -o serverHTTP-CGI
./serverHTTP-CGI

# Simula l'inserimento di numeri decimali e negativi (test)
curl -v "http://localhost:8000/sommatrice?op1=5.5&op2=-2.2"
\end{lstlisting}
\end{tcolorbox}

\end{document}
